\documentclass[a4paper]{article}

% Basic packages
% Español (México) con XeLaTeX: fontspec para fuentes Unicode, polyglossia
% para la división silábica y los nombres del documento en la variante mexicana.
\usepackage{fontspec}
\usepackage{polyglossia}
\setdefaultlanguage[variant=mexican]{spanish}
% Los nombres chinos van como «pinyin (original)»: xeCJK da a los caracteres
% Han su propia fuente; sin ella XeLaTeX los omite con un aviso.
% FandolSong no cubre algunos caracteres de nombres (U+73FA); WenQuanYi Zen Hei sí.
\usepackage[AutoFallBack=true]{xeCJK}
\setCJKmainfont{FandolSong-Regular.otf}
\setCJKfallbackfamilyfont{\CJKrmdefault}{WenQuanYi Zen Hei}
% polyglossia no traduce los nombres de listings.
\AtBeginDocument{\providecommand{\lstlistingname}{}\renewcommand{\lstlistingname}{Listado}%
  \providecommand{\lstlistlistingname}{}\renewcommand{\lstlistlistingname}{Índice de listados}}
\usepackage{amsmath, amssymb}
\usepackage{graphicx}
\usepackage[margin=2.5cm]{geometry}
\usepackage[most]{tcolorbox}
\usepackage{etoolbox}
\usepackage{listings}
\usepackage{hyperref}
\usepackage{booktabs}
\usepackage{subcaption}
\usepackage{float}

% Highlight boxes
\newtcolorbox{knowledgebox}[1]{
    enhanced, colback=blue!5!white, colframe=blue!75!black, colbacktitle=blue!75!black,
    coltitle=white, fonttitle=\bfseries, title={#1}, attach boxed title to top left={yshift=-2mm, xshift=2mm},
    boxrule=1pt, sharp corners
}

\newtcolorbox{importantbox}[1]{
    enhanced, colback=yellow!10!white, colframe=yellow!80!black, colbacktitle=yellow!80!black,
    coltitle=black, fonttitle=\bfseries, title={#1}, sharp corners
}

\newtcolorbox{warningbox}[1]{
    enhanced, colback=red!5!white, colframe=red!75!black, colbacktitle=red!75!black,
    coltitle=white, fonttitle=\bfseries, title={#1}, sharp corners
}

% Listings
\lstset{
    language=Python,
    basicstyle=\ttfamily\small,
    keywordstyle=\color{blue},
    stringstyle=\color{red!60!black},
    commentstyle=\color{green!60!black},
    breaklines=true,
    frame=single,
    numbers=left,
    numberstyle=\tiny\color{gray},
    captionpos=b,
    extendedchars=true
}

% Front-page metadata
\newcommand{\notetitle}{Notas de entrevista: Chen Tianqi (陈天奇) --- Sistemas de aprendizaje automático, largoplacismo y propósito inicial}
\newcommand{\noteauthors}{Elaboradas a partir de materiales públicos del curso}
\newcommand{\notedate}{2025-09-12}
\newcommand{\videochannel}{WhynotTV}
\newcommand{\videopublishdate}{2025-09-12}
\newcommand{\videoduration}{160 minutos}
\newcommand{\videourl}{https://www.bilibili.com/video/BV1s6pgzLE3y/}
\newcommand{\videocoverpath}{cover.jpg}
\newcommand{\repourl}{https://github.com/hqhq1025/ai-course-notes}

% Footer with repo info
\usepackage{fancyhdr}
\pagestyle{fancy}
\fancyhf{}
\fancyfoot[C]{\thepage}
\fancyfoot[R]{\footnotesize\color{gray}\href{\repourl}{AI Course Notes}}
\renewcommand{\headrulewidth}{0pt}
\renewcommand{\footrulewidth}{0pt}

\begin{document}

\begin{titlepage}
\centering
{\Large Notas de entrevista\par}
\vspace{1.2cm}
{\huge\bfseries \notetitle\par}
\vspace{0.8cm}
{\large \noteauthors\par}
\vspace{0.3cm}
{\large \notedate\par}
\vspace{1.2cm}

\ifdefempty{\videocoverpath}{
{\small Al generar las notas, indique la ruta local de la portada del video.\par}
}{
\includegraphics[width=0.82\textwidth,height=0.45\textheight,keepaspectratio]{\videocoverpath}\par
}

\vfill
\begin{tcolorbox}[width=0.9\textwidth, colback=black!2!white, colframe=black!60, sharp corners]
\textbf{Autor o canal del video}: \videochannel\par
\textbf{Fecha de publicación}: \videopublishdate\par
\textbf{Duración del video}: \videoduration\par
\ifdefempty{\videourl}{}{
\textbf{Enlace del video}: \href{\videourl}{\nolinkurl{\videourl}}\par
}\par
\textbf{Página del proyecto}: \href{\repourl}{\nolinkurl{\repourl}}\par
\end{tcolorbox}
\end{titlepage}

\tableofcontents
\newpage

%% ================================================================
%%  Inicio del cuerpo del texto
%% ================================================================

\section{Presentación del invitado y contexto de la entrevista}

El invitado de esta edición del WhynotTV Podcast (la tercera) es Chen Tianqi (陈天奇, Tianqi Chen), profesor asistente del departamento de aprendizaje automático de CMU y cofundador de OctoML (adquirida por NVIDIA a fines de 2024). Durante casi 20 años, Chen Tianqi ha impulsado, en torno al eje de \textbf{``hacer el aprendizaje automático más ligero, más rápido y más fácil de desplegar''}, proyectos de código abierto ampliamente adoptados como XGBoost, MXNet, TVM y MLC LLM. El artículo de XGBoost, en particular, tiene cerca de 70,000 citas y sigue siendo la herramienta más usada en el ámbito de los datos tabulares.

En 2019 publicó en Zhihu el artículo «Diez años de investigación en aprendizaje automático», que recibió decenas de miles de «me gusta». Seis años después, en 2025, ya lleva casi 20 años en este camino del aprendizaje automático. El conductor (WhynotTV) señala: ``Para mí, Chen Tianqi no es solo un nombre reconocido en el ámbito de la IA; también marcó mi primera comprensión del aprendizaje automático y de la investigación científica.''

Esta entrevista dura 160 minutos y parte de su infancia, pasa por la clase ACM de licenciatura, el laboratorio APEX de posgrado, los cuatro grandes proyectos de su doctorado en UW (XGBoost, MXNet, TVM y Net2Net durante su estancia en Google Brain), hasta llegar a su puesto docente en CMU y su emprendimiento con OctoML, para finalmente elevarse a una filosofía de vida sobre el largo plazo, la intención original y el valor.

\begin{importantbox}{El eje técnico de Chen Tianqi}
De XGBoost (2014) a MXNet (2015), a TVM (2017) y luego a MLC LLM (2023), cada proyecto de Chen Tianqi gira en torno al mismo problema central: \textbf{cómo lograr que los sistemas de aprendizaje automático funcionen de manera eficiente en distintos tipos de hardware, con menos recursos de ingeniería y de forma más automatizada}. Este eje atraviesa la era del big data, la era del aprendizaje profundo y la era de los modelos grandes. No es un investigador atado a un método en particular, sino alguien que persigue un problema------los métodos pueden cambiar (de los modelos de árboles a los marcos de aprendizaje profundo y a los compiladores), pero el problema sigue siendo el mismo: ``mejorar los sistemas de ML''.
\end{importantbox}

\section{Trayectoria de crecimiento: del condado a la clase ACM}

\subsection{La infancia y el primer encuentro con la computadora}

La secundaria de Chen Tianqi (陈天奇) fue la Secundaria No. 2 de Songyang, en el condado de Songyang, ciudad de Lishui, provincia de Zhejiang------su madre daba clases de preparatoria ahí, «en ese entonces sentía que estudiar en mi propio lugar también estaba bastante bien». El presentador le preguntó qué le interesaba más en su infancia, y dijo que la experiencia más interesante seguía siendo cómo empezó a tener contacto con la computadora.

La motivation inicial para aprender computación fue \textbf{querer hacer sus propios videojuegos}------«una motivation para aprender computación es poder crear juegos uno mismo. Claro que esa motivation nunca se llegó a cumplir». Pero justo ese deseo de «crear» lo llevó al mundo de la programación, y siguió presente después en su investigación------«hacer aprendizaje automático es lo mismo, crear inteligencia. Crear cosas es una motivation primigenia».

\subsection{Aprendizaje autodidacta de programación y competencias de informática}

En la secundaria básica, movido por su interés en hacer páginas web, buscó al profesor He (何), quien fue su guía inicial en la computación, y consiguió un libro introductorio de lenguaje C. En una época sin entrenadores de competencias ni GitHub, aprendió por su cuenta programación y competencias de informática a través de los Online Judge (OJ) en internet y de foros (como «Dalongshu» [大龙树]).

\begin{knowledgebox}{La ventaja del «camino poco convencional» del autodidacta}
En aquel entonces el lenguaje dominante en las competencias era Pascal, pero como Chen Tianqi solo había aprendido C al iniciarse, siguió compitiendo en C. Ese «camino poco convencional» terminó formando en él el hábito de no seguir las reglas establecidas. En segundo de preparatoria, por puro interés, dedicó un verano completo a escribir desde cero un transpiler (una versión primitiva de compilador) de Pascal a C. Esa experiencia le dio la actitud de «no tener miedo de hacer las cosas», y sentó las bases del estilo con el que después resolvería problemas partiendo de first principle.
\end{knowledgebox}

En segundo de preparatoria participó en la eliminatoria del NOIP y no pasó, pero asistió como oyente a la final y obtuvo el segundo lugar------en ese momento, ver el resultado «sí le dio bastante tristeza». Pero al año siguiente, en tercero de preparatoria, obtuvo el primer lugar, y además no descuidó en absoluto el gaokao, por lo que entró de forma normal a la Facultad de Electrónica e Información de la Universidad Jiao Tong de Shanghái. El presentador insistió: «¿en ese entonces siempre fuiste el primer lugar de todo el grado en tu condado?». Chen Tianqi respondió entre risas: «más o menos sí».

Vale la pena notar que la motivación inicial para aprender a programar en realidad fue querer hacer sus propios videojuegos------esa motivation nunca se llegó a cumplir, pero lo llevó a un mundo mucho más amplio.

\begin{warningbox}{Sobre «un ternero recién nacido no le teme al tigre»}
«Un ternero recién nacido no le teme al tigre» es una expresión que Chen Tianqi usa una y otra vez en la entrevista. Desde el estudio autodidacta de competencias en la preparatoria, pasando por escribir el transpiler, hasta después hacer MXNet para desafiar a TensorFlow y hacer TVM para entrar al campo de los compiladores, cada vez se lanzaba directamente sin saber «con qué problemas se iba a topar». Esta actitud se formó por haber aprendido por su cuenta en un condado pequeño, sin ninguna autoridad que le dijera «esto es muy difícil»------\textbf{al no haber visto nunca el techo, tampoco le teme al techo}.
\end{warningbox}

\subsection{La experiencia en la clase ACM de Jiao Tong}

En 2006, Chen Tianqi entró a la Facultad de Electrónica e Información de la Universidad Jiao Tong de Shanghái a través del gaokao, y luego, en una sesión informativa, conoció al profesor Yu Yong (俞勇); tras preparar sus materiales con mucho cuidado, se integró a la clase ACM. Recuerda que el profesor Yu (俞) le dijo después: «eres el estudiante que preparó sus materiales con más cuidado»------lo cual se parece mucho al personal statement que después prepararía para solicitar el PhD, «pero en ese momento no me di cuenta de eso». Como un «estudiante de gaokao venido de un lugar del que nadie había oído hablar», poder integrarse a la clase ACM fue «algo muy afortunado».

La influencia de la clase ACM sobre él se refleja en varias dimensiones:

\begin{enumerate}
    \item \textbf{Capacidad de expresión}: el «Foro Jiaozi» (教子讲坛) de la clase ACM exigía que cada estudiante eligiera un topic cada semestre y lo expusiera frente al grupo, lo cual le dio la confianza para volunteer directamente a dar presentaciones en las reuniones de grupo durante el doctorado.
    \item \textbf{«Ser humilde como persona, audaz en las acciones»}: un valor que el profesor Yu Yong enfatizaba una y otra vez. «Cuando se hace algo, hay que llevarlo al extremo» se convirtió en el motor central de su trabajo posterior en XGBoost.
    \item \textbf{El proyecto grande de compiladores}: el proyecto de diseño del curso de principios de compiladores en la clase ACM exigía construir un compilador desde cero; fue diseñado espontáneamente por los estudiantes de la primera generación, tomando como referencia la experiencia de universidades norteamericanas, y se fue transmitiendo de generación en generación. Chen Tianqi considera que este entrenamiento «sigue siendo de primer nivel incluso comparado con los estándares actuales de las universidades norteamericanas».
    \item \textbf{El proyecto grande de arquitectura de sistemas}: en la clase ACM también asumió el reto de ejecutar, en un CPU que él mismo había construido, el compilador que él mismo había escrito.
\end{enumerate}

\begin{importantbox}{La acumulación de bases sólidas en sistemas}
La experiencia del proyecto grande de compilador + CPU en la clase ACM sentó una base sólida para que Chen Tianqi diseñara después el compilador TVM y el acelerador VTA. Al recordar el desarrollo de TVM, dijo: «gracias a haber tenido esa experiencia anterior, no me sentí muy nervioso».
\end{importantbox}

\begin{importantbox}{El legado de «ser humilde como persona, audaz en las acciones»}
«Ser humilde como persona, audaz en las acciones», lo que el profesor Yu Yong enfatizaba una y otra vez en la clase ACM, influyó profundamente en Chen Tianqi. Después, en su carrera profesional, al tratar con distintas personas exitosas, descubrió que quienes de verdad logran las cosas tienen dos rasgos: primero, ser ambitious, tener el ímpetu para desafiar el statu quo (challenge the status quo); segundo, ser humble, saber que los factores del éxito «casi nunca vienen por completo de uno mismo», y que se necesita trabajo en equipo. La combinación de ambos------mantener el ímpetu dentro de la humildad------se convirtió en su estilo constante de actuar.
\end{importantbox}

\subsection{Resumen de la sección}

El crecimiento temprano de Chen Tianqi muestra un patrón que se repite: \textbf{aprender por su cuenta en un entorno con pocos recursos, no temerle al fracaso, resolver problemas partiendo de first principle}. Desde aprender OI por su cuenta en el condado, hasta escribir el transpiler en la preparatoria, hasta escribir en la clase ACM un compilador que se ejecutaba en su propio CPU, cada una de estas experiencias fue acumulando, para los grandes proyectos posteriores, el valor de «no tener miedo de hacer las cosas» y la intuición para el diseño de sistemas.
\section{Investigación temprana: el valor del fracaso}

\subsection{El laboratorio APEX y el primer contacto con el aprendizaje profundo}

Durante la licenciatura, Chen Tianqi (陈天奇) se unió al laboratorio APEX de la Universidad Jiao Tong, y con su compañero mayor Dai Wenyuan (戴文渊) (quien después fundó Fourth Paradigm, 第四范式) se inició en el aprendizaje automático leyendo el libro de texto clásico de Tom Mitchell. El motivo para elegir el aprendizaje automático fue sencillo: «que una máquina pueda aprender suena genial». Tom Mitchell fue precisamente el primer jefe del departamento de aprendizaje automático de CMU, y cuando Chen Tianqi (陈天奇) después llegó a dar clases en CMU, comentó con cierta emoción: «muchos profesores de aprendizaje automático crecieron leyendo ese libro».

En el tercer año, el profesor Xue (薛) les recomendó prestar atención al aprendizaje profundo, y él, junto con su compañero mayor Lin Yuan (林源), usó una máquina de Boltzmann restringida (RBM) para hacer clasificación en ImageNet. Para poder ejecutar los cálculos en GPU, compraron una tarjeta gráfica GTX 380 y descubrieron que la potencia de la fuente de poder del gabinete no era suficiente, así que pusieron una fuente de poder adicional fuera del gabinete, y así construyeron el primer servidor con GPU del laboratorio. Esto ocurrió en 2009--2010, dos años antes de AlexNet. En ese entonces, el grupo de Andrew Ng ya tenía artículos que discutían acelerar el aprendizaje profundo con GPU, y la herramienta CUDA ConvNet de Alex Krizhevsky ya se podía encontrar en línea, pero el aprendizaje profundo con GPU todavía estaba lejos de ser un consenso.

El proyecto de tesis de licenciatura de Chen Tianqi (陈天奇) fue precisamente sobre aprendizaje profundo, y además fue «escribir backpropagation desde cero en C++ y CUDA C»------en ese entonces, aunque ya existían las versiones previas de Theano y Caffe, él eligió implementarlo from scratch, «porque sentía que escribir muchas cosas desde cero todavía le permitía saber cómo se controla la forma en que se modela tu modelo». Construir un modelo tomaba seis meses de escribir código, y después venía ejecutar los experimentos.

\begin{warningbox}{Problema correcto + método incorrecto = dos años sin resultados}
Desde su proyecto de tesis de licenciatura hasta el segundo año de la maestría, Chen Tianqi (陈天奇) dedicó casi dos años a hacer «aprendizaje profundo on ImageNet», usando el método de RBM convolucional. El resultado nunca logró superar el baseline de SVM. Visto en retrospectiva, la dirección del aprendizaje profundo era correcta, pero el camino de RBM en sí no era el correcto------lo que finalmente work fue el supervised learning (como la red convolucional de AlexNet). Él lo resumió así: \textbf{no se puede fijar al mismo tiempo el problema que se quiere resolver y el método que se va a usar}. Una vez elegido el problema correcto, el método debe ajustarse con flexibilidad.
\end{warningbox}

\subsection{El giro decisivo en el KDD Cup}

En el último año de la maestría, el equipo participó en la competencia de sistemas de recomendación del KDD Cup. En ese momento sabían de un artículo llamado «Restricted Boltzmann Machines for Collaborative Filtering» (trabajo de Ruslan Salakhutdinov y otros), así que intentaron usar RBM en el sistema de recomendación. Pero pronto se dieron cuenta de que la factorización de matrices era más efectiva, y tras cambiar de método sin dudarlo obtuvieron buenos resultados. El proyecto de código abierto SVD Feature que Chen Tianqi (陈天奇) hizo después fue producto de ese trabajo; Danny (postdoc de Carlos Guestrin) descubrió a Chen Tianqi (陈天奇) a través de ese proyecto, le hizo una breve entrevista en su blog, y lo recomendó con Carlos y Alex Smola. Esta oportunidad llevó directamente a que obtuviera la offer de doctorado (PhD offer) de la UW.

También recibió al mismo tiempo una offer de CMU, pero como quería hacer investigación con Carlos, eligió la UW.

\begin{knowledgebox}{El «retorno implícito» de un proyecto de código abierto}
SVD Feature, como un proyecto temprano del que «ahora ya nadie sabe», se convirtió en la llave de entrada de Chen Tianqi (陈天奇) al círculo académico norteamericano. El valor de un proyecto de código abierto no está solo en el código en sí, sino más en que permite que tu trabajo sea descubierto por posibles colaboradores. Este patrón se repitió una y otra vez en cada uno de sus proyectos posteriores: el efecto de comunidad de XGBoost, la atención de los fabricantes de hardware que atrajo TVM, y la influencia de MLC LLM en la comunidad de dispositivos locales.
\end{knowledgebox}

\subsection{La pasantía en Hong Kong: la apertura de la visión de investigación}

Después de que el profesor Li Hang (李航) fue a dar una plática en APEX, Chen Tianqi (陈天奇) descubrió que «resulta que en el mundo existen todos estos otros problemas», y que, gracias a la acumulación de métodos que ya tenía, podía pensar en soluciones para muchos de ellos. Durante su pasantía en Hong Kong con los profesores Li Hang (李航) y Yang Qiang (杨强), empezó a entender que: \textbf{la investigación no consiste solo en resolver problemas, sino que, más importante aún, en saber qué problemas se pueden resolver}.

\begin{knowledgebox}{Dos estilos de investigación}
Chen Tianqi (陈天奇) se dio cuenta de que la investigación tiene dos style: (1) partir del método, convertirse en experto de cierto método y aplicarlo a distintos problemas; (2) partir del problema, interesarse en cierto tipo de problema y resolverlo con distintos métodos. Él pertenece al segundo tipo------\textbf{un investigador impulsado por el problema}. Esta es también la razón por la que pudo pasar de XGBoost a MXNet, y después a TVM, porque lo que él sigue no es un método en particular, sino el problema de «hacer que los sistemas de ML sean mejores».
\end{knowledgebox}

\subsection{Resumen de la sección}

La experiencia de «fracaso» en la investigación temprana trajo tres aprendizajes clave: (1) aprendió a aceptar el fracaso, lo que le dio más valor para actuar con soltura; (2) entendió que no se puede fijar al mismo tiempo el problema y el método; (3) mediante un entrenamiento básico extenso, acumuló «fuerza interna», lo que sentó las bases para entrar rápido a nuevos campos más adelante. Como él mismo dijo: \textbf{«si en ese momento hubiera tenido mucho éxito, en cambio podría haber generado dependencia de la trayectoria»}.
\section{XGBoost: de desafiar una hipótesis al estándar de la industria}

\subsection{Contexto de origen}

Después de que AlexNet irrumpiera en 2013, Chen Tianqi (陈天奇) discutió con su asesor Carlos Guestrin: ¿la razón del éxito del aprendizaje profundo era la red neuronal en sí misma, o un hypothesis space grande más muchos datos? Si esto último era la clave, entonces los modelos de árboles también tenían la capacidad de arbitrary grow (aumentar continuamente el número y la profundidad de los árboles), y tal vez también podrían lograr un avance con grandes volúmenes de datos.

Decidieron usar modelos de árboles para desafiar esta hipótesis. Basándose en su experiencia en KDD Cup, Chen Tianqi insistió en elegir el tree-based method (en lugar del kernel method que prefería su asesor). Consideraba que los modelos de árboles tenían una característica muy fuerte: podían arbitrary grow depth conforme aumentaba la cantidad de datos, mientras que el kernel method, aunque también era nonparametric, no tenía la misma ventaja en tabular data. Carlos le dio total libertad------``Estoy muy agradecido con mi asesor; me dio mucha libertad.''

El nombre de XGBoost es eXtreme Gradient Boosting; ``extreme'' refleja la búsqueda de llevar las cosas al extremo.

\begin{importantbox}{Las tres claves del éxito de XGBoost}
\begin{enumerate}
    \item \textbf{Rendimiento extremo}: al momento de su lanzamiento era la implementación de gradient boosting más rápida del mundo. Chen Tianqi llevó cada optimización al extremo, una tradición que proviene de la clase ACM: ``hacer una cosa y llevarla al extremo''.
    \item \textbf{Comunidad de código abierto}: el éxito de XGBoost proviene en gran medida de las contribuciones de la comunidad. Por ejemplo, un abogado francés contribuyó con un paquete de visualización en R, usado para detectar fraude fiscal. La fuerza de la comunidad llevó el proyecto mucho más allá de lo que una sola persona podría lograr.
    \item \textbf{Enfoque}: hacer un solo algoritmo, y hacer ese algoritmo bien, llevándolo al extremo. En comparación con el proyecto anterior ``grande y abarcador'' SVD Feature, la estrategia de enfoque de XGBoost era más adecuada para un equipo pequeño.
\end{enumerate}
\end{importantbox}

\subsection{Integración vertical entre algoritmo y sistema}

XGBoost incorpora un manejo automático de missing value: si el valor de alguna característica falta, el modelo decide automáticamente cómo clasificarlo con base en el pattern histórico, sin que el usuario necesite hacer una limpieza de datos previa. Esta función ``era muy, muy popular'', porque en los datos reales los valores faltantes están en todas partes. El método tradicional consiste en rellenarlos con la mediana u otros medios, pero XGBoost deja que el propio modelo aprenda la manera óptima de manejarlos.

Al trabajar en XGBoost, Chen Tianqi se dio cuenta por primera vez de que: \textbf{para construir un sistema de aprendizaje automático, no basta con hacer solo el algoritmo o solo el sistema; es necesario integrarlos verticalmente}. Este entendimiento atravesó todos sus proyectos posteriores. Señaló que el éxito de XGBoost no se debía solamente a que corría rápido a nivel de sistema (multihilo, out-of-core learning, cómputo distribuido), sino también a la optimización detallada a nivel de algoritmo (regularización, sparsity-aware split finding, missing value handling). \textbf{Solo entendiendo al mismo tiempo las necesidades del algoritmo y las restricciones del sistema se puede construir una herramienta verdaderamente útil}.

\begin{knowledgebox}{La filosofía de escritura del artículo de XGBoost}
El artículo de XGBoost no se escribió hasta el tercer año del PhD. La exigencia de Carlos no era conseguir la aceptación del artículo, sino ``escribir un artículo del que el lector pueda aprender algo''. Ese mismo año, en KDD, el laboratorio de Carlos presentó otro artículo, LIME (Explainable ML), realizado por Marco Ribeiro, que dio origen a esta línea de investigación: el aprendizaje automático interpretable. El style de investigación de Carlos consistía en alentar a los estudiantes a realizar trabajos capaces de fundar un campo nuevo, en lugar de hacer mejoras incrementales sobre direcciones ya existentes.
\end{knowledgebox}

\subsection{Reflexión sobre los modelos de árboles frente a las redes neuronales}

Al mirar en retrospectiva más de una década después, considera que el juicio de aquel momento fue ``partly right, partly wrong'':

\begin{itemize}
    \item \textbf{La parte correcta}: la capacity del modelo no debería estar limitada. Los modelos de árboles pueden arbitrary grow depth, una característica muy fuerte. En tabular data y time series, los modelos de árboles siguen siendo hasta hoy el go-to tool.
    \item \textbf{Lo que no se anticipó}: las redes neuronales lograron la expansión de capacity mediante una forma ``más bruta''------no fue grow el parameter space de manera gradual, sino establecer un parameter space enorme desde el principio.
    \item \textbf{La propiedad combinatoria de las redes neuronales}: módulos como el transformer block y el linear block pueden combinarse entre sí, de modo que la capacidad expresiva se representa de forma más compact mediante la combinación.
\end{itemize}

\subsection{La influencia de su asesor Carlos Guestrin}

Carlos exigía muchísimo a sus estudiantes: en el primer encuentro ya les decía ``You should write the best paper in the world''. Su laboratorio tenía la norma de ``no beamer, only PowerPoint'', para insistir en que se usaran medios visuales que permitieran al lector comprender rápidamente el contenido de la investigación. El artículo de XGBoost no se escribió hasta el tercer año del doctorado; la exigencia de Carlos era ``escribir un artículo del que el lector pueda aprender algo'', no optimizar la tasa de aceptación.

\subsection{Resumen de la sección}

El éxito de XGBoost no fue casualidad: un rendimiento llevado al extremo, una comunidad activa y una estrategia de enfoque, sumados a la integración vertical entre algoritmo y sistema, lo convirtieron en el estándar de la industria, con casi 70,000 citas. Pero en ese momento Chen Tianqi no anticipó este éxito------simplemente ``hicimos algo que nos pareció decente, y luego lo liberamos como código abierto''.
\section{MXNet: el florecimiento de los marcos de aprendizaje profundo}

\subsection{De CXXNet a MXNet}

El éxito de AlexNet reavivó el entusiasmo de Chen Tianqi (陈天奇) por el aprendizaje profundo. Él y Xu Bin (许斌) (ahora en NVIDIA) desarrollaron CXXNet, que usaba la técnica de programación con plantillas de C++ (expression template) para generar kernels de GPU en línea. La idea central de esta técnica es: cuando escribes la update rule (por ejemplo, operaciones de multiplicación, resta, etc.), las plantillas de C++ pueden comprimir esas operaciones en un solo kernel y generar código de GPU en línea. CXXNet tenía un buen desempeño y también se usó para competir en concursos de Kaggle relacionados con aprendizaje profundo.

Al mismo tiempo, el campo de los marcos de aprendizaje profundo estaba en una etapa de «florecimiento pleno»:

\begin{itemize}
    \item \textbf{Zhang Zhen (张真) y Minjie (敏捷)} (MSR/NYU): desarrollaron Minerva, cuya filosofía de diseño Python first era muy adelantada para su época
    \item \textbf{Li Mu (李沐)} (CMU): trabajó en Parameter Server (servidor de parámetros), que necesitaba que lo distribuido fuera first class citizen
    \item \textbf{Zhang Chiyuan (张池元)} (MIT): escribió en Julia un marco llamado Moka
    \item \textbf{Lin Min (林敏)} (NUS): desarrolló el marco Purine
    \item \textbf{Yang Qing (杨青)} (UC Berkeley): Caffe ya había tenido su release y era el marco dominante de la época, pero su núcleo era C++
\end{itemize}

Gracias a la oportunidad que ofrecían congresos académicos como NeurIPS, estos estudiantes de doctorado dispersos entre UW, CMU, MIT, NUS y NYU se reunieron y decidieron unir fuerzas para construir un marco unificado------este es el origen de MXNet.

\begin{knowledgebox}{El origen del nombre de MXNet}
MXNet es una mezcla (Mix) de Minerva y CXXNet, y al mismo tiempo representa «mix programming style»------la combinación de grafos de cómputo declarativos (declarative, donde primero se declara el grafo de cómputo y después se ejecuta) y de programación imperativa (imperative, donde se calcula en cuanto se escribe). En ese momento existía un trade-off evidente entre estos dos paradigmas: lo declarativo favorecía la optimización del sistema (permitía optimizaciones globales), mientras que lo imperativo era más amigable para el desarrollador (el debug era más simple). MXNet optó por dar soporte a ambos. Más tarde PyTorch eligió el modo puramente imperativo (eager mode), TensorFlow 1.x eligió el modo puramente declarativo (aunque la versión 2.0 después también migró a eager mode), y MXNet ofreció desde el inicio un modo híbrido.
\end{knowledgebox}

Chen Tianqi también mencionó que entre sus compañeros de laboratorio en UW estaba Joe Redmon (el creador de YOLO), «él ya había hecho YOLO antes de su doctorado». Después de la aparición de AlexNet, muchas personas del laboratorio comenzaron a explorar la dirección del aprendizaje profundo.

\subsection{Filosofía central de diseño}

Las tres grandes innovaciones de MXNet:
\begin{enumerate}
    \item \textbf{Python First}: en la época en que Caffe tenía a C++ como núcleo, MXNet eligió a Python como ciudadano de primera clase.
    \item \textbf{Planificación automática}: a partir de las dependencias de cómputo, implementaba de forma automática el paralelismo entre múltiples GPU, de modo que la comunicación y el cómputo pudieran hacer overlap.
    \item \textbf{Diferenciación automática}: evolucionó de la backpropagation escrita a mano en Caffe hacia una diferenciación automática basada en el grafo de cómputo.
\end{enumerate}

\subsection{La competencia con TensorFlow}

Durante un verano en que hacía una pasantía en Google Brain, Chen Tianqi descubrió que TensorFlow estaba a punto de lanzarse. Él fue uno de los usuarios dogfood de TensorFlow, y al mismo tiempo compartió con el equipo de Brain la experiencia de diseño de MXNet. Pero el equipo no se echó atrás: «un ternero recién nacido no le teme al tigre, nosotros también podíamos compete un poco». Incluso siguieron «avanzando con mucha pasión».

En la línea temporal, el proyecto de MXNet se inició antes del lanzamiento de TensorFlow. Todos se habían reunido en el NeurIPS del año anterior para decidir seguir adelante; TensorFlow tuvo su release a finales de 2,015, y el lanzamiento oficial de MXNet fue aproximadamente en la misma época. Al final, MXNet solo envió un workshop paper, con los autores ordenados alfabéticamente------«sencillamente sentíamos que hacer un proyecto juntos era algo bueno».

\begin{knowledgebox}{La colaboración «pura» entre estudiantes de doctorado}
El presentador preguntó: con tantos doctorandos reunidos y sin que nadie estuviera a cargo, ¿cómo podía work out? Chen Tianqi respondió: «en realidad no había nadie que tuviera la última palabra. Porque lo que todos querían, sencillamente, era hacer el proyecto juntos». Este modelo de colaboración basado en intereses compartidos y sin relaciones jerárquicas tenía, en cambio, sus ventajas dentro de la comunidad académica de código abierto------cada quien contribuía por pasión técnica y no por KPI. Dentro de la colaboración, «hablar de la técnica en términos puramente técnicos es en realidad muy simple: se discute la propuesta directamente y se avanza».
\end{knowledgebox}

\begin{warningbox}{Las lecciones del ocaso de MXNet}
MXNet terminó siendo superado por PyTorch, y Chen Tianqi resumió dos lecciones clave:
\begin{enumerate}
    \item \textbf{La experiencia de usuario tiene prioridad sobre el desempeño}: durante mucho tiempo, MXNet fue el mejor marco en paralelismo entre múltiples GPU (NVIDIA lo usaba para ejecutar MLPerf), pero en experiencia de usuario no estaba a la altura de PyTorch. «A la hora de elegir entre experiencia de usuario y desempeño, pensábamos que había que tener ambos, pero en realidad primero había que elegir la experiencia de usuario.»
    \item \textbf{Falta de personal de ingeniería}: dar soporte a cada nueva generación de GPU exigía una gran inversión de ingeniería, y un equipo de estudiantes de doctorado difícilmente podía sostenerlo. Esto lo llevó directamente a orientarse hacia TVM------\textbf{usar tecnología de compilación para reducir la carga de ingeniería}.
\end{enumerate}
\end{warningbox}

\subsection{El papel de pilar de Li Mu (李沐)}

Chen Tianqi mencionó especialmente que Li Mu (李沐) fue quien hizo «un esfuerzo y un sacrificio enormes» en el proyecto MXNet. «En realidad, al final, quien más impulsó todo esto fue Li Mu.» Después de graduarse, Li Mu se fue a Amazon, impulsó que MXNet se convirtiera en el marco oficial de AWS y desarrolló GluonCV y GluonNLP. Chen Tianqi señaló que GluonNLP «in some sense se parece un poco a Hugging Face»------ofrecía una interfaz unificada para distintos modelos preentrenados. Pero, al final, estos esfuerzos a nivel de ecosistema no lograron cambiar el panorama en el que la comunidad de PyTorch ya se había consolidado.

Un hecho poco conocido es que, durante mucho tiempo, NVIDIA estuvo haciendo contribute a MXNet y usándolo. La razón es que, entre todos los marcos de código abierto, el desempeño de MXNet en paralelismo entre múltiples GPU siempre fue el mejor. Al enviar sus resultados de MLPerf, NVIDIA «cuando de verdad necesitaba obtener el mejor puntaje, siempre usaba MXNet»------y solo fue quedando atrás poco a poco después de que ResNet dejó de ser el protagonista del benchmark.

Chen Tianqi también mencionó que Wang Naiyan (王乃岩) y Hou Xiaodi (侯晓迪), de TuSimple (图森未来), «durante mucho tiempo estuvieron ayudando a contribuir y a dar soporte a MXNet». Tanto el éxito de este proyecto como su ocaso final no se pueden entender sin las contribuciones de muchas personas de la comunidad.

\subsection{Resumen de la sección}

La experiencia con MXNet le enseñó dos cosas a Chen Tianqi: (1) la experiencia de usuario importa más que los indicadores técnicos; (2) el cuello de botella de personal de ingeniería debe resolverse con medios técnicos más fundamentales (un compilador), y no simplemente sumando más gente. Estas dos lecciones dieron origen directamente a TVM.
\section{TVM: la apertura de la compilación de aprendizaje automático}

\subsection{La motivación que parte de un problema}

De CXXNet a MXNet, Chen Tianqi (陈天奇) descubrió un cuello de botella que se repetía: cada generación de GPU y de hardware nuevo requería una gran cantidad de kernels escritos a mano, que debían admitir distintos layouts de datos (como NHWC frente a NCHW), distintos tamaños de entrada y distintas formas de hardware; el espacio de combinaciones era enorme. Los kernels escritos a mano «sencillamente no bastaban».

\begin{importantbox}{El problema central que TVM debía resolver}
¿Se podía usar la técnica de compilación para generar automáticamente estos kernels? ¿Se podía lograr que menos personas construyeran rápidamente, con la técnica correcta, sistemas de aprendizaje automático orientados a hardware nuevo? Este era el punto de partida de TVM------\textbf{que la complejidad de ingeniería de los sistemas de aprendizaje automático dejara de crecer linealmente con la cantidad de tipos de hardware}.
\end{importantbox}

\subsection{La diferencia entre GPU, TPU y NPU, y su convergencia}

En la entrevista, Chen Tianqi explicó la diferencia entre tres tipos de chips de cómputo:

\begin{itemize}
    \item \textbf{GPU}: proviene de la programación gráfica; su característica central es SIMT (Single Instruction Multiple Thread), que programa con una gran cantidad de hilos en paralelo. Con el desarrollo de CUDA, ya «no tiene relación» con la programación gráfica.
    \item \textbf{TPU/NPU}: cuenta con un conjunto de instrucciones más complejo, donde una sola instrucción puede completar una operación matricial masiva o una lectura/escritura de memoria a gran escala. Por lo general no tiene el concepto de hilo, porque una sola instrucción ya es, de por sí, muy «grande».
    \item \textbf{Tendencia de convergencia}: las GPU incorporaron el Tensor Core (que originalmente solo tenían las NPU), y las NPU también empezaron a incorporar el concepto de hilo para acercarse al ecosistema de CUDA. «Ahora este límite es cada vez más difuso.»
\end{itemize}

Esta tendencia de convergencia es precisamente la manifestación del valor persistente de TVM------el modelo de programación de cada generación de hardware cambia, y se necesita la técnica de compilación para tender un puente entre las diferencias.

\subsection{El proceso de desarrollo: «construir una ciudad en una isla desierta»}

TVM se inspiró en parte en Halide (un trabajo de Jonathan Ragan-Kelley y otros, del MIT, que usaba la técnica de compilación para optimizar pipelines en el campo del procesamiento de imágenes), pero la compilación de aprendizaje automático era, en ese momento, un terreno completamente vacío. Chen Tianqi describió esta experiencia como «construir una ciudad en una isla desierta»------no había ningún camino ya trazado al que recurrir.

Desde que empezó a idearlo, al volver del GTC de abril de 2016, hasta la primera versión ejecutable, en marzo de 2017, pasaron cerca de 11 meses, la mayor parte dedicados a escribir código. Recuerda que este proceso fue de «iteración constante»: «Aunque no lo entendiera, no tenía miedo; era como un becerro recién nacido que no le teme al tigre.»

Mientras hacía TVM, junto con Thierry Moreau, su colaborador del área de arquitectura de computadoras, construyó con FPGA un acelerador de NPU programable. Su idea inicial era más «wild»------meter directamente la red neuronal, mediante bake, en el cableado (wire) del FPGA (más adelante alguien hizo un trabajo relacionado con HLS). Al final decidieron hacer primero una NPU parecida a una TPU, pero con una filosofía de diseño completamente distinta de la corriente dominante en ese momento.

\begin{warningbox}{El pago de la experiencia del trabajo grande de la clase ACM}
Al recordar el desarrollo de TVM, Chen Tianqi mencionó varias veces su experiencia en la clase ACM. «En ese entonces, en la clase ACM teníamos un trabajo grande de arquitectura de computadoras, y yo me propuse el reto de ejecutar, en un CPU que yo mismo había escrito, un compilador que yo mismo había escrito. TVM se parece muchísimo a esa experiencia------también era construir desde cero un compilador for un hardware nuevo. Gracias a esa experiencia anterior, sentí que no me ponía tan nervioso.» Esto muestra que un proyecto retador que, en la etapa de licenciatura, parece una «actividad extracurricular», puede, años después, dar un pago de una forma inesperada.
\end{warningbox}

Chen Tianqi señaló en particular que, en ese entonces, él «no venía formado en arquitectura de computadoras»; pensaba por completo from first principle cómo hacerlo. El conductor comentó: «Siempre has sido así------en el concurso de bachillerato eras un becerro recién nacido que no le teme al tigre, escribiendo el compilador en la clase ACM eras un becerro recién nacido que no le teme al tigre, haciendo XGBoost eras un becerro recién nacido que no le teme al tigre, haciendo TVM sigues siendo un becerro recién nacido que no le teme al tigre. ¿No es esta tu mayor ventaja?» Chen Tianqi respondió: «Creo que elegir el problema y cómo abordarlo es muy importante. Hacer un problema mediocre y hacer un problema importante toman casi el mismo tiempo, así que hay que elegir sin falta el problema correcto.»

Al mismo tiempo, junto con su colaborador del área de arquitectura de computadoras, construyó con FPGA un acelerador de NPU programable (VTA). El planteamiento de diseño de esta NPU tenía una visión muy adelantada a su tiempo:

\begin{knowledgebox}{El diseño innovador del acelerador VTA}
La mayoría de las NPU tempranas eran completamente hardwired (no programables, con una sola función fija). Pero VTA adoptó un diseño \textbf{programable a nivel de instrucción}, que permitía implementar múltiples operadores con un mismo conjunto de microinstrucciones. Para resolver el problema del traslape (overlap) entre la lectura/escritura de memoria y el cómputo, se inspiraron en el concepto de hyper-threading: dividieron lógicamente la NPU en dos hilos, uno dedicado a las operaciones de memoria y otro al cómputo, sincronizados mediante un mecanismo de señales de hardware. Este diseño «sigue siendo bastante interesante incluso hoy», y en la programación asíncrona más reciente de NVIDIA se pueden ver rastros similares.
\end{knowledgebox}

\subsection{El sistema de coordenadas de Jeff Dean y la elección de un tema de alto riesgo}

Durante su estancia como pasante en Google Brain, Jeff Dean llevó a Chen Tianqi a una sala de juntas y dibujó un sistema de coordenadas: un eje era el riesgo (risk), el otro el impacto (impact). Entre las distintas opciones de proyecto, Chen Tianqi eligió el tema con mayor riesgo y también mayor impacto------Net2Net. El problema de este proyecto era: dado un modelo pequeño ya entrenado, ¿cómo hacer un bootstrap rápido de un modelo grande, de modo que su capacidad pudiera aumentar?

La razón por la que Chen Tianqi eligió este tema fue que siempre tuvo una visión: la inteligencia necesita poder evolucionar, y la capacidad del modelo necesita poder crecer. Si una red neuronal empieza siendo pequeña, en cierto punto deja de ser suficiente. Desde la perspectiva del lifelong learning, ¿cómo resolver este problema? Señaló: «De hecho, ahora te das cuenta de que este problema todavía existe de la misma manera------esta es también la razón por la que ahora hay una tendencia hacia el scaling up a partir de modelos pequeños.»

Jeff Dean era, en ese entonces, su mentor en Brain, y además «tenía la mente muy abierta». Chen Tianqi fue, en ese momento, prácticamente el primer intern de Jeff Dean------en una etapa en la que «su GAN tal vez todavía no se había publicado, o apenas se acababa de publicar».

Su criterio de elección era: \textbf{«El tiempo es limitado, así que hay que hacer sin falta lo que uno quiere hacer. No tiene que ser algo con impacto, pero sí tiene que ser algo que uno quiera hacer.»}

Ante la pregunta de «si les recomendarías a los investigadores jóvenes elegir temas de alto riesgo y alto rendimiento», su respuesta fue muy pragmática: puede ser distinto para cada persona. Pero como él ya había pasado por muchos fracasos, «fracasar una vez más tampoco es tan grave». «Almost always terminamos eligiendo una dirección de riesgo relativamente alto y retorno relativamente alto.»

\subsection{El valor durante el posdoctorado y la historia de la Recomputation}

Antes de hacer TVM, Chen Tianqi había pensado en un método para «entrenar modelos más profundos con una complejidad menor que lineal» (ahorrando memoria mediante recomputation), pero Carlos le dijo que este trabajo, publicado tal cual, «sería apenas un NeurIPS poster», que no era suficientemente profundo, y que si se iba a hacer, había que llevarlo al nivel de best paper------había que diseñar un método más automatizado para elegir la estrategia de recomputation. Como su energía era limitada, no continuó por esa línea, y más adelante hubo un follow-up work que hizo un trabajo relacionado.

La exigencia de Carlos no era negar el valor de esa línea, sino decir: \textbf{«Tienes que elegir lo que haces, y una vez que decides hacerlo, tienes que hacerlo a fondo.»}

Li Mu (李沐), en un comentario en Zhihu, reveló que, en el tercer o cuarto año de su doctorado, Chen Tianqi le dijo que iba a hacer TVM, y Li Mu le advirtió que el riesgo era demasiado alto------abrir un pozo nuevo y grande cuando el doctorado estaba por terminar. La respuesta de Chen Tianqi fue: «Sin importar si en el futuro voy a trabajar, a emprender o a buscar una posición académica, lo más importante sigue siendo hacer, en ese momento, lo que uno mismo considera más importante.» Li Mu comentó después: «Visto ahora, esta es la única forma correcta de actuar.»

\begin{importantbox}{La lógica de fondo para tomar decisiones}
Al tomar decisiones, Chen Tianqi se hace una pregunta a sí mismo: \textbf{«Si elijo este camino y fracaso, ¿me voy a arrepentir?»} Si la respuesta es que no se arrepentiría, lo elige. Esto no es un análisis racional de costo-beneficio, sino una manera de asegurarse de seguir lo que siente en el corazón. «Muchas veces de lo que uno se arrepiente no es de si logró el éxito, sino de si tomó o no el camino que idealmente quería tomar.»
\end{importantbox}

\subsection{La evolución de TVM y su nicho ecológico}

Desde su nacimiento hasta ahora, TVM ha pasado por múltiples rondas de reestructuración, desde la compilación de grafos inicial hasta Relax (un compilador de grafos rediseñado para los LLM). Chen Tianqi enfatizó que la filosofía de diseño de TVM es \textbf{«evolucionable»}------reinventarse continuamente para adaptarse a las nuevas necesidades. Señaló: «Si comparas el TVM actual con su forma en el momento en que empezó, ya hemos pasado por varias generaciones de iteración.»

Frente a competidores como Torch Compile, XLA y Triton, considera que la necesidad de la técnica de compilación no solo no ha disminuido, sino que va en aumento:

\begin{itemize}
    \item \textbf{Especialización (specialized) del hardware}: la forma en que avanza cada generación de chips es distinta, y el modelo de programación cambia incluso entre generaciones sucesivas de un mismo vendor
    \item \textbf{El «entrelazamiento» entre el modelo y el hardware}: para alcanzar la eficiencia óptima, el diseño del modelo se acopla cada vez más con las características del hardware------DeepSeek, al diseñar en función de un entorno de hardware específico, es un ejemplo típico de esto
    \item \textbf{Hardware Lottery}: el hardware limita las posibilidades de desarrollo del modelo, y se necesita la técnica de compilación para romper esta restricción
    \item \textbf{El auge de los Agentes de IA}: la aparición de los sistemas de agentes y de los DSL trae nuevos escenarios de aplicación para la técnica de compilación
\end{itemize}

TVM tiene actualmente dos direcciones de desarrollo: una es integrarse mejor con ecosistemas como PyTorch (abstrayendo por separado una especificación general de llamadas a funciones, para facilitar las llamadas entre distintos frameworks); la otra es lograr avances en dominios verticales (como compilar y desplegar LLM en páginas web y en dispositivos vehiculares). Su meta última es: \textbf{«hacer el desarrollo de ingeniería de aprendizaje automático tan simple como escribir una red neuronal»}------algo que, por ahora, todavía está lejos de lograrse.

\begin{warningbox}{Reinventarse a uno mismo frente a hacer patch a un sistema existente}
Dentro de OctoML, Chen Tianqi también enfrentó la disyuntiva entre reinvent y patch. Para dar soporte a los LLM, echaron abajo la parte de compilación de grafos original de TVM y rehicieron Relax. Esta decisión «no la iba a aprobar todo el mundo», porque la técnica ya existente podía resolver el problema del momento a corto plazo. Pero él insiste en que: «If you don't reinvent yourself, there will be a time where you bite the technical debt.» A largo plazo, la acumulación de deuda técnica es más temible que el costo a corto plazo de una reestructuración.
\end{warningbox}

\subsection{Resumen de la sección}

TVM es la obra que reúne, de la manera más completa, el estilo de investigación de Chen Tianqi de «partir del problema». Surgió de un cuello de botella de ingeniería descubierto durante el desarrollo de MXNet, se inspiró en el campo de los compiladores, y en su diseño fusionó la idea del codiseño de hardware y software. Aunque enfrenta una competencia intensa, el lugar central de la técnica de compilación en los sistemas de aprendizaje automático ya es un consenso.
\section{MLC LLM: la implementación en el dispositivo en la era de los modelos grandes}

\subsection{De TVM a MLC LLM}

MLC LLM es una aplicación vertical de TVM------un motor de inferencia de LLM multiplataforma capaz de compilar e implementar modelos de lenguaje grandes en diversos dispositivos como iPhone, Android, tarjetas gráficas de AMD y MacBook. Utiliza tecnología de compilación de aprendizaje automático, porque el equipo «sabe con toda claridad que no tenemos la capacidad suficiente para extender manualmente el soporte a cada plataforma».

\subsection{El futuro de la IA en el dispositivo}

Chen Tianqi (陈天奇) compara el panorama actual de la IA con la era de las macrocomputadoras: «Todos usan mainframe y piensan que todo se puede centralize. Pero luego llegó la personal computer. Today la gente sigue acostumbrada a usar el software que tiene en su propia computer; aunque la vision de algo como Chromebook no está mal, la mayoría sigue prefiriendo tener su personal computer.» Los factores impulsores centrales de la implementación en el dispositivo incluyen:

\begin{enumerate}
    \item \textbf{Privacidad (Privacy)}: aunque muchas personas dicen que «no care» (el presentador bromeó diciendo: «si entras a mi ChatGPT y le preguntas por todos mis social security number, seguramente ya los sabe todos»), la necesidad de privacidad existe de manera objetiva
    \item \textbf{Tiempo real}: el presentador, desde la perspectiva del campo de la robótica, señaló que los robots necesitan tomar decisiones en el orden de los milisegundos------«si se va a caminar o a tomar una taza, no hay tiempo de quedarse ahí parado reason 10 segundos». NVIDIA está trabajando recientemente en integrar el chip Thor e implementarlo en robots, lo cual es precisamente un reflejo de esta necesidad
    \item \textbf{Costo}: cuando el preentrenamiento empieza a «start a bit» (el crecimiento de la capacidad del modelo se desacelera), la gente pasa de perseguir el «next level of intelligence» a centrarse en application y cost
    \item \textbf{Especialización}: una vez que el local model alcanza un nivel «suficiente» en un ámbito específico, el panorama cambiará------«no se trata de igualar por completo al modelo más grande en el ámbito general, sino que, si en algunos ámbitos especializados se puede llegar a un nivel suficiente, I think the picture can change»
\end{enumerate}

Chen Tianqi mismo también señaló que, si existiera una IA local con una capacidad casi equivalente a la de GPT, «I would opt to try that way». La variable clave es qué tan fuerte puede llegar a ser realmente la capacidad de los modelos en el dispositivo.

\begin{knowledgebox}{La analogía histórica de Mainframe a Personal Computer}
Esta analogía resulta muy reveladora. La historia de la computación pasó de las macrocomputadoras centralizadas a las computadoras personales distribuidas, y es posible que la IA repita esta trayectoria. Pero no se trata de una simple repetición------la nube y el dispositivo se repartirán las funciones según el escenario, tal como hoy coexisten las aplicaciones web y las aplicaciones locales. La variable central que determina el panorama es el límite de capacidad del modelo local, junto con la ponderación que hace el usuario entre costo, privacidad y tiempo real.
\end{knowledgebox}

\subsection{Resumen de la sección}

MLC LLM refleja la filosofía constante de Chen Tianqi: reducir la barrera de ingeniería mediante tecnología de compilación, para que incluso equipos más pequeños puedan implementar IA en distintos tipos de hardware. El futuro de la IA en el dispositivo depende del equilibrio entre la capacidad del modelo y el costo, pero es muy probable que la tendencia de largo plazo sea un panorama distribuido donde «florezcan cien flores».
\section{20 años de evolución de los sistemas de aprendizaje automático}

\subsection{Las tres grandes olas}

Chen Tianqi divide los últimos 20 años de evolución de los sistemas de ML en tres etapas, y en cada una de ellas la innovación en sistemas impulsó directamente los avances algorítmicos:

\textbf{Primera etapa: la era del big data (2008--2012)}

El evento emblemático fue la competencia de sistemas de recomendación Netflix Challenge, además de empresas como Jinri Toutiao que despegaron gracias a los sistemas de recomendación. El avance de esta época vino de que el volumen de datos saltó de unos cientos de puntos de datos a cientos de millones. Los sistemas representativos incluyen Apache Spark, XGBoost, y la regresión lineal a gran escala que Li Mu (李沐) y Dai Wenyuan (戴文渊) hicieron en Baidu (un sistema no abierto).

\textbf{Segunda etapa: el avance del aprendizaje profundo (2012--2017)}

El éxito de AlexNet no fue solo un triunfo del algoritmo, sino también un triunfo de la ingeniería de sistemas: Alex Krizhevsky entrenó el modelo durante una semana con dos GPU, una proeza en un momento en que los experimentos solían durar apenas unos minutos. La aparición de Caffe ``impulsó en muy buena medida que el campo de la visión por computadora adoptara el aprendizaje profundo''. Después, con la llegada de MXNet, TensorFlow y PyTorch, construir un modelo pasó de ``seis meses escribiendo 20,000 líneas de CUDA'' a ``unos minutos escribiendo unas cuantas líneas de Python''.

\textbf{Tercera etapa: la era de los modelos grandes (2020--presente)}

El transformer domina la arquitectura de los modelos, y la atención se ha desplazado hacia la ley de escalamiento, el entrenamiento distribuido y la inferencia eficiente. Los nuevos desafíos de sistemas incluyen: clústeres de entrenamiento del orden de decenas de miles de GPU, estrategias de paralelismo de modelos, el despliegue en el dispositivo del usuario final, y la generación estructurada (como XGrammar).

\begin{table}[H]
\centering
\begin{tabular}{p{2.5cm}p{3cm}p{3cm}p{4cm}}
\toprule
\textbf{Etapa} & \textbf{Periodo} & \textbf{Evento emblemático} & \textbf{Sistemas representativos} \\
\midrule
Era del big data & 2008--2012 & Netflix Challenge & Spark, XGBoost \\
Aprendizaje profundo & 2012--2017 & AlexNet, ImageNet & Caffe, MXNet, PyTorch \\
Era de los modelos grandes & 2020--presente & GPT, transformer & vLLM, Triton, MLC LLM \\
\bottomrule
\end{tabular}
\caption{Las tres grandes etapas de la evolución de los sistemas de aprendizaje automático}
\end{table}

\subsection{El papel habilitador de los sistemas frente a los algoritmos}

Chen Tianqi enfatiza que detrás de cada avance algorítmico hay una innovación a nivel de sistemas:

\begin{itemize}
    \item AlexNet pudo triunfar porque Alex Krizhevsky era un ``hacker muy fuerte'' que entrenó el modelo durante una semana con dos GPU, en un momento en que los experimentos solían durar apenas unos minutos
    \item Caffe ``impulsó en muy buena medida que el campo de la visión por computadora adoptara el aprendizaje profundo''
    \item La aparición de los marcos de aprendizaje profundo hizo que construir un modelo pasara de ``seis meses escribiendo 20,000 líneas de CUDA'' a ``unos minutos escribiendo unas cuantas líneas de Python''
\end{itemize}

\begin{importantbox}{La relación triangular entre algoritmo, datos y cómputo}
``Ahora todo el mundo sabe que la inteligencia artificial depende de los datos, el cómputo y el algoritmo. Pero cuando empezamos, casi toda la atención estaba puesta en el algoritmo.'' Fue el progreso de los sistemas (la aceleración por GPU, el entrenamiento distribuido, los marcos de diferenciación automática) lo que permitió a los investigadores de algoritmos iterar sus experimentos con rapidez, una relación habilitadora que suele pasarse por alto.
\end{importantbox}

\subsection{Oportunidades y desafíos en el ámbito académico}

La brecha de recursos que enfrenta la investigación en sistemas de ML dentro de la academia es real: un clúster de cien GPU en la universidad frente a las decenas de miles de GPU de los laboratorios de vanguardia. El presentador preguntó directamente: ``Entre cien y diez mil GPU todavía hay una diferencia de dos o tres órdenes de magnitud.'' Pero Chen Tianqi considera que esta brecha no es tan determinante como se piensa:

\begin{itemize}
    \item \textbf{La escasez de recursos estimula la creatividad}: ``Lack of resource is also an opportunity to make people more creative.'' Por eso él sigue dedicándose a la compilación de aprendizaje automático: en la universidad es imposible resolver el problema simplemente sumando 100 personas.
    \item \textbf{Colaboración con la comunidad de código abierto}: la universidad no puede actuar como un individuo aislado; necesita producir trabajo que la comunidad de código abierto pueda construir en conjunto. Proyectos exitosos como VRL y Magiline tienen componentes universitarios.
    \item \textbf{Avances verticales}: no hace falta hacerlo todo, pero sí se puede llevar un módulo clave al máximo nivel. Por ejemplo, XGrammar (generación de datos estructurados), que desarrollaron ellos mismos, ya es uno de los estándares de facto de la comunidad de código abierto para la generación estructurada: que la salida JSON de un agente tenga el formato correcto es un problema ``muy importante, pero en realidad ninguna de las soluciones actuales lo resuelve bien''.
    \item \textbf{La brecha a nivel de investigación y desarrollo no es tan grande}: ``Muchas veces la industria tampoco es tan rica, sobre todo si no eres un equipo que en verdad está haciendo pre-training.'' El Flame Center, fundado recientemente en CMU, cuenta con un clúster del orden de cien GPU, suficiente para hacer mucha optimización de kernels e investigación de sistemas.
\end{itemize}

\begin{warningbox}{Sobre la ansiedad de que ``ya se resolvieron todos los problemas''}
Chen Tianqi contó una anécdota: cuando empezaba en el aprendizaje automático, Dai Wenyuan estaba sentado a su lado escuchando una presentación sobre optimización de SVM, y suspiró: ``si hubiera entrado a este campo dos años antes, habría alcanzado todos esos problemas elegantes, que parece que ya se resolvieron por completo''. Eso fue hace 20 años. ``It turns out there is always new problem.'' Cada época tiene sus propios problemas y desafíos nuevos; lo importante es ser capaz de descubrirlos y resolverlos.
\end{warningbox}

\subsection{El rumbo futuro de MLSYS}

Chen Tianqi ha sido durante mucho tiempo un impulsor importante detrás de la MLSys Conference. Considera que los sistemas de aprendizaje automático son, ``igual que los sistemas de bases de datos, una dirección guiada por problemas''. Las direcciones importantes de los próximos años incluyen:

\begin{enumerate}
    \item \textbf{Optimización del entrenamiento e inferencia de modelos grandes}: cómo hacer más eficientes el pre-training, el fine-tuning y el serving
    \item \textbf{Adaptación de hardware} (compilación de aprendizaje automático): la puesta en marcha del hardware y su adaptación en software siguen siendo un enorme cuello de botella
    \item \textbf{AI for Systems}: usar agentes de IA para optimizar la propia ingeniería de sistemas de ML
    \item \textbf{Infraestructura para sistemas de agentes}: generación estructurada (como XGrammar), tool calling, entre otros
\end{enumerate}

\begin{knowledgebox}{¿La convergencia de los modelos hará decaer a MLSYS?}
El presentador planteó una pregunta incisiva: los modelos de lenguaje de gran escala son cada vez más convergentes (los cinco modelos principales podrían concentrar el 99\% del volumen de inferencia); ¿no le quitaría esto su razón de ser a MLSYS? Chen Tianqi reconoce que esta posibilidad no puede descartarse por completo, pero la situación real es más compleja. Por un lado, el hardware mismo ``has to be specialized'': el modelo de programación cambia en cada generación de chips. Por otro lado, el ``Hardware Lottery'' implica que la forma actual de los modelos quizá no sea la óptima: el florecimiento del aprendizaje profundo se debió justamente a que coexistieron cien enfoques distintos. MLSYS no consiste solo en optimizar los pocos modelos actuales, sino en habilitar la exploración de nuevas formas de modelo.
\end{knowledgebox}

\subsection{Resumen de la sección}

La evolución de los sistemas de aprendizaje automático siempre ha ido de la mano con los avances algorítmicos. Los sistemas no son un simple adorno, sino tecnología habilitadora. La academia conserva ventajas propias en este campo: neutralidad, la paciencia para invertir a largo plazo, y una afinidad natural con la comunidad de código abierto. Chen Tianqi subraya en particular que, aunque la academia no puede competir en escala con los laboratorios de vanguardia, ``la escasez de recursos estimula la creatividad'': llevar un módulo clave al máximo nivel y ampliar su influencia mediante la colaboración con la comunidad de código abierto es el camino propio de la academia.
\section{OctoML: de código abierto a startup}

\subsection{Motivación para emprender}

La fundación de OctoML fue muy natural: varios estudiantes y profesores del proyecto TVM vieron una oportunidad de comercialización y decidieron hacer una startup. La motivación central de Chen Tianqi (陈天奇) era «que la comunidad de código abierto pudiera tener éxito»: hacer una startup por el código abierto, no hacer código abierto por la startup.

Durante el primer año participó de tiempo completo, y después pasó a part-time, encargándose principalmente de la dirección técnica y de la investigación de la próxima generación de tecnología. Él tenía «muy claro que no quería ser CEO», porque «tienes que renunciar a mucha tecnología», y su passion estaba en hacer buena tecnología.

\subsection{El cambio de modelo de negocio}

El modelo de negocio de OctoML pasó por un cambio significativo:

\begin{enumerate}
    \item \textbf{Etapa temprana (2019--2022)}: mediante la compilación de aprendizaje automático, ayudaban a los fabricantes de hardware y a los clientes a desplegar modelos con mayor rapidez en distintos tipos de hardware. Los clientes potenciales incluían fabricantes de chips y empresas que necesitaban desplegar ML.
    \item \textbf{Etapa posterior (2023--2024)}: se transformaron en un proveedor de servicios de API de LLM (similar a Fireworks y Together), ofreciendo directamente endpoints de modelos como Llama, con cobro por llamada a la API. La tecnología de compilación de aprendizaje automático pasó a segundo plano, convirtiéndose en tecnología de soporte que ayudaba a reducir el costo de inference.
    \item \textbf{Finales de 2024}: fueron adquiridos por NVIDIA, y el equipo quedó a cargo de la dirección del AI compiler de NVIDIA.
\end{enumerate}

El presentador insistió en que el mercado de API tenía una competencia intensa, en la que todos competían por tener el Llama inference cost más bajo, y preguntó cómo respondía OctoML ante eso. Chen Tianqi reconoció que la competencia era intensa, pero señaló que aún había espacio para diferenciarse en la forma del producto, por ejemplo, dando soporte a multi-LoRA inference. También admitió con franqueza que la adquisición por NVIDIA no fue una ruta de salida planeada de antemano: «Nunca pensé en esa pregunta; en ese momento solo sentía que bastaba con hacer bien una cosa y ver qué pasaba.»

\begin{knowledgebox}{El precedente de la experiencia emprendedora del mentor}
Antes de unirse a OctoML, Chen Tianqi ya había sido testigo de la evolución de Turi (anteriormente GraphLab/Dato), la empresa fundada por su mentor Carlos Guestrin: partió como software de procesamiento de grafos a gran escala en CMU y, para cuando fue adquirida, ya se había transformado por completo en una plataforma de Data Science. «Así que, al empezar a hacer esto, ya tenía la preparación mental de que lo que empiezas a hacer y what you end up doing sin duda serán distintos.» Databricks partió de Spark, y ahora Spark tampoco es ya su parte más importante. Esta es una regularidad general de las startups de código abierto.
\end{knowledgebox}

\begin{warningbox}{La relación sutil entre el código abierto y el negocio}
Chen Tianqi admitió que hacer una startup a partir de un proyecto de código abierto enfrenta el problema de que «el código abierto es tan bueno que nadie paga». Pero su postura es: (1) el código abierto es tu punto de partida, pero casi nunca es tu punto de llegada; el negocio central de Databricks tampoco es ya Spark; (2) lo que toda startup piensa hacer al inicio y lo que termina haciendo «sin duda serán distintos»; (3) una vez que eliges el código abierto, hay que aceptar tanto sus ventajas (retroalimentación de la comunidad, influencia) como sus retos (el diseño del modelo de negocio).
\end{warningbox}

\subsection{El Lesson más grande del emprendimiento}

Cinco años de emprendimiento le dieron a Chen Tianqi importantes conocimientos más allá de lo técnico:

\begin{itemize}
    \item \textbf{Human factor is very important}: «El lesson más grande de hacer una startup es que muchas veces no solo hay factores técnicos. Cómo comunicarte con la gente, cómo desempeñar un role de leadership dentro de un equipo grande: creo que esto es sumamente importante.»
    \item \textbf{Hay que reinvent yourself}: no puedes aferrarte sin cambios a tu ruta técnica original. «Si hace dos años hubiéramos seguido optimizando ResNet, creo que al final habríamos perdido de forma muy contundente.» En el campo de la IA, que cambia con rapidez, la velocidad y el valor para transformarse técnicamente son fundamentales.
    \item \textbf{La excellence técnica nunca es un error}: «Esta opinión no la comparte todo el mundo. Pero, desde mi punto de vista, muchas veces hacer tecnología excelente nunca es un error. Porque in the end, you have to have something to differentiate yourself.»
    \item \textbf{Idealismo + realismo}: «Tengo un poco de idealismo. Pero esta experiencia me hace sentir que soy más realista: sé cómo, por medios realistas, lograr mejor mis ideales.»
    \item \textbf{Product Market Fit}: hay que pensar a fondo a quién y qué problema resuelve realmente tu producto. Pero, al mismo tiempo, no puedes fijarte solo en el PMF de corto plazo e ignorar el valor técnico de largo plazo.
    \item \textbf{Tener claro tu propio rol}: «Tengo muy claro que no quiero ser CEO. Porque ser CEO es sumamente difícil, tienes que renunciar a mucha tecnología. Mi passion está en hacer buena tecnología.» Conocer en qué eres bueno y qué amas, y luego encontrar al co-founder adecuado que te complemente.
\end{itemize}

\begin{warningbox}{La relación sutil entre el código abierto y el negocio}
Chen Tianqi admitió que hacer una startup a partir de un proyecto de código abierto enfrenta el problema de que «el código abierto es tan bueno que nadie paga» (el presentador puso el ejemplo de Anyscale y Ray). Pero su postura es: (1) el código abierto es tu punto de partida, pero casi nunca es tu punto de llegada; el negocio central de Databricks tampoco es ya Spark; (2) lo que toda startup piensa hacer al inicio y lo que termina haciendo «sin duda serán distintos»; (3) una vez que eliges el código abierto, hay que aceptar tanto sus ventajas (retroalimentación de la comunidad, influencia, confianza de los usuarios) como sus retos (el modelo de negocio necesita diferenciarse). Hu Yuanming (胡渊鸣), invitado del episodio anterior, con su Taichi (太极) también enfrentó un problema similar: «Mucha gente dice que Taichi es muy bueno, pero no está dispuesta a pagar.» Este es un reto general para el emprendimiento basado en código abierto.
\end{warningbox}

\subsection{Resumen de la sección}

Los cinco años de trayectoria de OctoML muestran la ruta típica de comercialización de un proyecto académico de código abierto: de estar impulsado por la tecnología a estar impulsado por el producto, de herramienta de compilación a servicio de API, hasta terminar en una adquisición estratégica. En este proceso, Chen Tianqi aprendió dimensiones importantes más allá de lo técnico: el factor humano, el ajuste entre producto y mercado, y la necesidad de una autorrenovación constante.
\section{La convicción y la práctica del código abierto}

\subsection{Por qué insiste en el código abierto}

Chen Tianqi (陈天奇) es un firme creyente del open source y la open science. Considera que el valor del código abierto se manifiesta en varios niveles:

\begin{enumerate}
    \item \textbf{Intercambio comunitario}: compartir mutuamente, conocer a personas distintas, obtener la retroalimentación positiva de que le den «me gusta» a tu código
    \item \textbf{Una vía para que la universidad llegue a la práctica}: a la universidad misma no le resulta fácil llevar la tecnología directamente a la práctica; el código abierto es la mejor manera de compartirla con todos y construirla en conjunto
    \item \textbf{Formación de estudiantes}: en una empresa, uno normalmente empieza por el componente pequeño de un proyecto grande, pero un doctorante que hace un proyecto de código abierto tiene la oportunidad de \textbf{diseñar la arquitectura de un sistema desde cero}------ese tipo de entrenamiento es sumamente valioso
    \item \textbf{Una ventaja natural de la universidad}: neutralidad más espíritu de entrega a la comunidad
\end{enumerate}

\subsection{La clave para construir una comunidad de código abierto}

\begin{importantbox}{Los tres elementos de un proyecto de código abierto exitoso}
\begin{enumerate}
    \item \textbf{La tecnología correcta}: que todos believe in it. La tecnología misma es la base; sin una buena tecnología, la comunidad no se congrega.
    \item \textbf{La passion de mantenerlo de manera continua}: el éxito de un software no lo determina la forma que tenga su código en un momento dado, sino que necesita mantenerse constantemente. Los issues y la retroalimentación de los usuarios son «sumamente valiosos»------son el mejor insumo para la mejora iterativa.
    \item \textbf{Build trust}: si uno hace algo y de repente deja de mantenerlo, cuando lance el siguiente proyecto la gente dudará de «si en verdad hay que creer en esto». La confianza requiere acumularse a largo plazo y, una vez perdida, es muy difícil reconstruirla.
\end{enumerate}
\end{importantbox}

Chen Tianqi mismo suele «revisar con frecuencia» los issues y la retroalimentación de los usuarios. Considera que el código abierto no es free------requiere mucho tiempo para revisar la retroalimentación, corregir errores e interactuar con la comunidad, cosas que en el sistema tradicional de evaluación de la investigación no tienen retribución directa (no cuentan como publicaciones). «Así que no es que todos los grupos deban tomar este camino; casi ningún grupo lo hará, pero es una manera de hacerlo.»

El presentador preguntó de nuevo: ¿revisas uno por uno todos los issues de todos tus proyectos de código abierto? Chen Tianqi respondió que sí los revisa con frecuencia, porque «un software no puede tener su éxito determinado por la forma que tenga su código en un momento dado------necesita mantenerse constantemente».

\subsection{Resumen de la sección}

Para Chen Tianqi, el código abierto no es solo una forma de difundir tecnología, sino también un valor: que la tecnología la use más gente, que cada quien pueda crear su propio sistema de IA. Esta también fue la intención original detrás de OctoML------«para que la comunidad de código abierto pueda prosperar».
\section{Carrera académica: el profesor que sigue escribiendo código en primera línea}

\subsection{Elegir CMU}

Los tres motivos por los que Chen Tianqi eligió la carrera académica: le gusta la investigación (disfruta la incertidumbre), le gusta resolver problemas nuevos junto con estudiantes, le gusta enseñar a la gente a construir cosas. Da clases en el departamento de aprendizaje automático de CMU y a la vez dirige a estudiantes de doctorado en investigación de sistemas de ML. Ya cerca de graduarse había decidido seguir el camino académico en lugar de dedicarse por completo a la industria.

Cuando le preguntaron «¿cómo convences a la gente más talentosa de hacer un doctorado en lugar de irse a XAI u OpenAI a trabajar en infraestructura?», respondió: «Lo que distingue al doctorado de otros caminos es que it's like a longer investment (es como una inversión más larga): tienes suficiente tiempo y suficiente libertad para hacer ciertas cosas.» En la academia, la alineación con la open science y la escala de tiempo de tolerancia al error siempre son mayores. Cree que «siempre habrá estudiantes con ideas dispuestos a elegir este camino».

\subsection{Insistir en escribir código en primera línea}

A diferencia de muchos profesores que «desde el primer día como profesores nunca volvieron a escribir código», Chen Tianqi hasta hoy sigue haciendo git commit en primera línea. En el último año ha estado reconstruyendo desde cero el núcleo de TVM, «la parte de foundation (fundamento) la hago yo». Escribir código es para él algo «recreational» (recreativo), más parecido a una afición.

\begin{knowledgebox}{Las herramientas de IA hacen posible que un profesor siga escribiendo código}
Chen Tianqi lo reconoce: «Si no fuera por las AI based tools (herramientas basadas en IA), hacer esta reconstrucción me resultaría mucho más difícil. Pero with the helpful tools (con estas herramientas de ayuda) (él usa Cursor), it's possible (es posible).» Las herramientas de IA para programar bajan el umbral para mantenerse hands-on (con las manos en la masa), y le permiten conservar el contacto con la línea técnica incluso mientras dirige a su equipo. Él considera: «Muchas veces necesitas estar en primera línea, para get a good sense of what's going on (tener una buena intuición de lo que está pasando).»
\end{knowledgebox}

\subsection{La metáfora de pescar: acumulación y momento oportuno}

Ante la pregunta de «si el ritmo de la investigación en IA se ha acelerado, ¿todavía tiene sentido hacer proyectos a escala de años?», Chen Tianqi usó una metáfora:

\textbf{«Investigar es como pescar. Puedes tomar un tridente y ensartar al pez cuando pasa nadando, pero muchas veces fallas el golpe. Otra manera es tejer poco a poco una red, y cuando llegan los peces, después de lanzarla varias veces tal vez logres atraparlos.»}

Él considera que todo trabajo sobresaliente requiere acumulación; muchas veces lo que parece un resultado «lanzado rápidamente» en realidad es fruto de una acumulación de largo plazo.

\subsection{Estilo de asesoría}

Su advising style (estilo de asesoría) continúa la tradición de Carlos y añade su propio sello:

\begin{enumerate}
    \item \textbf{Hacer trabajo con impacto}: esto viene heredado de la frase de Carlos «write the best paper on the wall» (escribe el mejor artículo de la pared)
    \item \textbf{Dar importancia a la presentación}: continúa la tradición de Carlos de «no beamer, only PowerPoint» (nada de Beamer, solo PowerPoint), con énfasis en que el artículo y la exposición permitan a la gente aprender rápido lo que hiciste
    \item \textbf{Fomentar el trabajo en equipo}: en el área de sistemas de ML, llevar un proyecto de la investigación a una adopción real requiere colaboración en equipo
    \item \textbf{Formar un independent student} (estudiante independiente): sostiene firmemente que el estudiante debe tener la capacidad de drive things forward (impulsar las cosas hacia adelante), mientras él mismo act as a support role (funge como apoyo)
    \item \textbf{Hacer pocos proyectos y profundizar en cada uno}: «deliberately (deliberadamente) queremos hacer menos cosas, y concentrarnos en los proyectos que consideramos que se pueden profundizar»; esta es también la razón por la que, si solo se mira el número de artículos publicados, su producción no parece «eficiente»
\end{enumerate}

Su manera de administrar el tiempo consiste en elegir y enfocarse: saber qué es lo que más quiere y luego hacerlo. Al mismo tiempo, tiene la fortuna de contar con estudiantes muy independent (independientes) y con un co-founder (cofundador) que comparten la operación diaria.

\subsection{Resumen de la sección de la carrera académica}

Chen Tianqi eligió un camino poco común: profesor, más emprendedor, más alguien que escribe código en primera línea. Administra su tiempo mediante el enfoque (hacer pocos proyectos y profundizar en cada uno), la delegación (la fortuna de tener estudiantes independent y un co-founder) y la ayuda de las herramientas de IA. Con la metáfora de «pescar» subraya la importancia de la acumulación: en vez de perseguir tendencias con un tridente, es mejor tejer poco a poco una red en el propio campo.
\section{Filosofía de vida: largoplacismo, espíritu inicial y valentía}

\subsection{Los fundamentos matemáticos del largoplacismo}

Chen Tianqi (陈天奇) escribió en «Estos diez años de investigación en aprendizaje automático»: «Muchas de las dificultades y del desamparo no son más que parte de una fluctuación aleatoria; basta con invertir suficiente tiempo y paciencia para que el proceso aleatorio termine convergiendo a un estado estable acorde con lo invertido.»

Dice que esta frase «es literalmente un teorema»: el time average de un proceso aleatorio en efecto converge a una distribución estacionaria. Pero una cosa es conocer el teorema y otra muy distinta es perseverar en medio de los reveses.

\begin{importantbox}{La capacidad de aceptar el fracaso}
«El fracaso tampoco es tan terrible; después hay que aceptar el hecho de que uno puede fracasar.» Esta convicción se ha ido reforzando desde que en las Olimpiadas de Informática (OI) de la preparatoria solo obtuvo el segundo lugar, pasando por dos años sin resultados haciendo aprendizaje profundo durante la licenciatura, hasta los diversos giros de su emprendimiento. Chen Tianqi considera que: \textbf{mientras más fracasos se experimentan, más atrevimiento se tiene para actuar con soltura}, porque «fracasar tampoco es para tanto, de todos modos sigo viviendo bien». El fracaso temprano incluso es una especie de «protección»: si uno tiene demasiado éxito desde el principio, puede terminar generando dependencia de trayectoria.
\end{importantbox}

\subsection{La valentía proviene del yo del pasado}

Una idea profunda de Chen Tianqi es que \textbf{la valentía para actuar no aumenta automáticamente con la experiencia}. Por el contrario, mientras más cosas se ven, más fácil resulta volverse tímido y vacilante.

«Creo que muchas veces tu valentía tal vez sea el compromiso que tu yo del pasado hizo contigo mismo. Tienes que cumplir el compromiso que tu yo pasado hizo consigo mismo para poder seguir insistiendo en este style de trabajo.»

Él describe este proceso como un ciclo: el yo del pasado le da valentía al yo del presente, el yo del presente se anima continuamente a sí mismo durante el proceso de hacer las cosas, y luego le entrega esa valentía al yo del futuro.

\subsection{Mantener el espíritu inicial es lo más difícil}

Cuando le preguntaron «mirando hacia los próximos 5--10 años, ¿qué quieres hacer?», la respuesta de Chen Tianqi fue inesperada:

\textbf{«Creo que lo más difícil para mí ahora sigue siendo mantener el espíritu inicial.»}

Explica que, a medida que aumentan los recursos, aumentan las opciones y crece la responsabilidad, resulta más difícil mantener la actitud de aquellos años, la de «un becerro recién nacido no le teme al tigre». Cuando trabajaba en MXNet no le preocupaba en absoluto la competencia de TensorFlow, pero ahora, frente a un proyecto nuevo, no puede evitar preguntarse «si en realidad va a tener éxito o no». Mantener el espíritu inicial requiere un recordatorio deliberado hacia uno mismo.

El presentador insiste: los recursos que puedes movilizar son varios órdenes de magnitud mayores que hace 10 años, ¿por qué sientes que es más difícil mantener el espíritu inicial? Chen Tianqi responde: «Obtener recursos implica necesariamente responsabilidad. Tienes que garantizar que el equipo pueda tener éxito, tienes que producir el impact correspondiente para justify hacer las cosas relacionadas. Si te dejas arrastrar por estas presiones, pierdes la valentía para hacer cosas risky.»

Su estrategia para enfrentarlo es: seguir haciendo cosas más risky dentro del equipo de CMU, «estar dispuesto a renunciar y estar dispuesto a aceptar el fracaso, esta actitud en realidad es la misma que al principio». Los recursos y el equipo han crecido, y en efecto se pueden hacer más cosas que antes, pero el cambio de actitud puede hacer que uno «olvide esa actitud propia».

\begin{warningbox}{La trampa mental después del éxito}
Esta discusión revela una paradoja profunda: el éxito y la acumulación de experiencia pueden \textbf{erosionar la valentía necesaria para innovar}. Mientras más cosas se ven y más riesgos se conocen, más fácil resulta volverse conservador. Chen Tianqi considera que la manera de contrarrestar esta tendencia es volver al espíritu inicial: no tomar la decisión mediante un análisis racional de costos y beneficios, sino preguntarse a uno mismo: «¿es este el camino que idealmente quiero recorrer?». Si la respuesta es sí, entonces hacerlo sin importar el resultado.
\end{warningbox}

\subsection{Sobre el éxito y la felicidad}

En la parte final, el presentador hizo varias preguntas más personales.

\textbf{Sobre el éxito}: Chen Tianqi considera que el éxito tiene muchas definiciones: influencia, riqueza, avances tecnológicos, etcétera. Pero para él, lo más importante es «hacer cosas con la conciencia tranquila». «Poder influir en la gente en el ámbito técnico y poder contribuir en conjunto al mundo ya es algo bueno.»

\textbf{Sobre el mensaje para su yo futuro}: si le pidieran decirle una frase a su yo del pasado y otra a su yo del futuro, su elección no sería regresar a cambiar algo, sino todo lo contrario: \textbf{necesita que su yo del pasado le recuerde a su yo presente y futuro}: recordar el compromiso consigo mismo, sostener el ideal y seguir adelante.

\textbf{Sobre el estado ideal de vida}: el Chen Tianqi ideal de dentro de 10 o 20 años debería ser «un Chen Tianqi que ha tenido más fracasos, y que puede seguir manteniendo esa actitud inicial para desafiar lo que queremos desafiar».

\subsection{Resumen de la sección sobre la filosofía de vida}

La filosofía de vida de Chen Tianqi puede resumirse en tres palabras clave: \textbf{largoplacismo} (creer que el proceso aleatorio converge), \textbf{espíritu inicial} (hacer cosas con la conciencia tranquila) y \textbf{valentía} (aceptar el fracaso, sin dejarse arrastrar por los recursos y la responsabilidad). Estas no son consignas vacías, sino principios de acción verificados una y otra vez a lo largo de 20 años de investigación y emprendimiento. Vale la pena señalar que él no considera que estas cualidades se refuercen naturalmente con el aumento de la experiencia; por el contrario, considera que mantener el espíritu inicial y la valentía requiere un recordatorio constante y deliberado hacia uno mismo.
\section{Síntesis y ampliación}

\subsection{Un hilo conductor técnico claro}

Al revisar los casi 20 años de carrera profesional de Chen Tianqi (陈天奇), un hilo conductor lo atraviesa por completo: \textbf{hacer que los sistemas de aprendizaje automático sean más eficientes, más fáciles de usar y más fáciles de desplegar}. Cada proyecto fue una respuesta al cuello de botella del proyecto anterior:

\begin{itemize}
    \item Al participar en el KDD Cup descubrió que necesitaba modelos de árboles eficientes $\rightarrow$ \textbf{XGBoost}
    \item Al trabajar en XGBoost reavivó su entusiasmo por el aprendizaje profundo $\rightarrow$ \textbf{CXXNet/MXNet}
    \item Al trabajar en MXNet descubrió que la mano de obra de ingeniería era el cuello de botella $\rightarrow$ \textbf{TVM} (sustituir los kernels escritos a mano por técnicas de compilación)
    \item TVM necesitaba validarse en un dominio vertical $\rightarrow$ \textbf{MLC LLM} (despliegue de modelos grandes en dispositivos de borde)
    \item La necesidad de comercializar TVM $\rightarrow$ \textbf{OctoML} (código abierto $\rightarrow$ empresa emergente $\rightarrow$ adquirida por NVIDIA)
\end{itemize}

\begin{table}[H]
\centering
\begin{tabular}{p{2.2cm}p{1.8cm}p{9cm}}
\toprule
\textbf{Proyecto} & \textbf{Periodo} & \textbf{Problema que resolvió} \\
\midrule
XGBoost & 2014 & Permitió que los modelos de árboles se ejecutaran más rápido y fueran más fáciles de usar en datos a gran escala (manejo automático de missing value) \\
MXNet & 2015 & Permitió que el marco de trabajo de aprendizaje profundo diera soporte a Python first, programación automática, paralelismo con múltiples GPU y diferenciación automática \\
TVM & 2017 & Permitió que los modelos de ML se adaptaran automáticamente a distintos tipos de hardware (GPU/TPU/NPU), reduciendo los kernels escritos a mano \\
MLC LLM & 2023 & Permitió que los modelos de lenguaje grandes se ejecutaran en dispositivos de borde como teléfonos, navegadores y MacBook \\
\bottomrule
\end{tabular}
\caption{Cronología de las contribuciones técnicas de Chen Tianqi}
\end{table}

La lógica subyacente de este hilo conductor es: \textbf{los sistemas de ML no deberían ser tan complejos; deberían permitir que menos personas los construyeran rápidamente con la técnica correcta}. Esta convicción se mantiene constante desde la era de XGBoost (una sola persona llevándolo al extremo), pasando por la era de TVM (el compilador genera automáticamente los kernels), hasta la era de MLC LLM (despliegue multiplataforma en dispositivos de borde).

\subsection{Lecciones para los investigadores jóvenes}

Con base en el contenido de la entrevista, los consejos centrales que Chen Tianqi da a los jóvenes pueden resumirse en los siguientes puntos, cada uno respaldado por su experiencia personal:

\begin{enumerate}
    \item \textbf{Elegir el problema correcto}: «Trabajar en un problema mediocre y en uno importante toma casi el mismo tiempo.» Esto se lo enseñó el profesor Ao Ping (敖平) de la Universidad Jiao Tong, y lo ha confirmado una y otra vez a lo largo de 20 años. ImageNet fue el problema correcto y las RBM fueron el método equivocado: la elección del problema importa más que la elección del método.

    \item \textbf{No fijar al mismo tiempo el problema y el método}: esta es la «lección más valiosa» que aprendió tras dos años sin resultados trabajando en ImageNet en la Universidad Jiao Tong. Una vez elegido el problema, el método debe ajustarse con flexibilidad: más adelante, el cambio decidido de las RBM a la factorización de matrices en el KDD Cup fue la aplicación de esta lección.

    \item \textbf{Llevarlo al extremo}: la exigencia de Carlos era «write the best paper in the world». No se trataba de perseguir la cantidad de artículos, sino de llevar cada proyecto al nivel de best paper. El éxito de XGBoost es el mejor ejemplo de «llevar un algoritmo al extremo».

    \item \textbf{Aceptar el fracaso}: «Trabajar en un problema muy difícil implica necesariamente incertidumbre y fracaso.» Fracasar temprano es, en realidad, algo bueno: te vuelve más capaz de actuar con libertad y evita la dependencia de la trayectoria.

    \item \textbf{Seguir el instinto}: al elegir, pregúntate «si fracaso, ¿me arrepentiré?». No es un análisis racional de costo-beneficio, sino una manera de asegurarte de que sigues el camino ideal.

    \item \textbf{Partir de first principle}: no temer entrar en campos desconocidos. Desde escribir un transpiler hasta construir un compilador y diseñar un NPU, Chen Tianqi siempre se lanzó de inmediato, como «un ternero recién nacido que no le teme al tigre». La capacidad de diseñar sistemas se afina con la práctica.

    \item \textbf{El código abierto es una oportunidad única para los estudiantes universitarios}: en una empresa empiezas por un componente pequeño dentro de un proyecto grande; en un proyecto de código abierto tienes la oportunidad de diseñar la arquitectura de un sistema desde cero. Este tipo de entrenamiento es difícil de conseguir en otro lugar.
\end{enumerate}

\begin{knowledgebox}{Un consejo especial para los estudiantes de doctorado}
Cuando le preguntaron: «si hoy fueras un estudiante de doctorado de tercer año y solo pudieras elegir corregir o cambiar una cosa, ¿qué harías?», la respuesta de Chen Tianqi fue: «Creo que elegiría seguir haciendo lo que hago ahora.» Es decir, continuar con la compilación de aprendizaje automático: permitir que menos personas, con menos engineering resource, resuelvan el problema del bring up de los sistemas de ML en el hardware más reciente. No cambiaría de rumbo solo porque la época es distinta, porque este problema «está lejos de resolverse».
\end{knowledgebox}

\subsection{Perspectivas sobre el futuro de la IA}

A Chen Tianqi no le interesa definir qué es la AGI («no me interesa demasiado cómo se defina la AGI en sí»); le importa más que sea «pragmatically útil»: ¿puede un equipo de una sola persona generar la producción de un equipo de 100 personas? La tecnología actual «sí nos permite avanzar más rápido hacia esa meta»: por ejemplo, él mismo usa Cursor para refactorizar TVM, y sin herramientas de IA «refactorizar sería mucho más difícil».

Respecto a cómo se configurará el mundo de la IA en el futuro, espera ver que florezcan cien flores: que el código abierto y los individuos también puedan innovar, en lugar de que unos cuantos gigantes dominen todo. Cuando le preguntaron si «la probabilidad de que aparezca ese mundo está aumentando o disminuyendo», su respuesta fue cautelosa: «Solo puedo decir que requiere el esfuerzo de cada persona. Sin duda, distintas personas tendrán creencias distintas.»

La línea técnica en la que insiste siempre ha sido: \textbf{hacer más con menos recursos y de manera más automática}. «Tal vez se resuelva en unos años, tal vez no se resuelva necesariamente, pero al menos es un problema muy interesante.»

Esta frase quizá sea el mejor resumen de toda la carrera de Chen Tianqi: no persigue el éxito garantizado, sino los problemas interesantes. Mientras un problema sea suficientemente importante e interesante, vale la pena invertir en él: incluso si el resultado es el fracaso, también es «un proceso de volver a partir».

\subsection{Conclusión}

Al final de la entrevista, el conductor dejó un espacio abierto para que Chen Tianqi hablara libremente. Dijo:

\textbf{«A lo largo de todo el proceso he tenido mucha suerte de recibir la ayuda de muchísimas personas. Lo que más profundamente me ha dejado es esto: todavía podemos aceptar el fracaso. En el proceso de aceptar el fracaso, hay que seguir animándose a hacer cosas interesantes. Cada proceso es, en realidad, un proceso de volver a partir. We need to reinvent ourselves, then keep doing the next thing.»}

Si tuviera que decirle una frase a su yo del pasado y otra a su yo del futuro, su elección no sería mirar hacia atrás desde el presente para enseñarle algo a su yo pasado, sino al contrario: necesita que aquel yo pasado, «un ternero recién nacido que no le teme al tigre», le recuerde a su yo presente: \textbf{recordar el compromiso consigo mismo, mantener el ideal y seguir adelante.}

El conductor resumió la entrevista diciendo: el Chen Tianqi ideal de dentro de 10 o 20 años debería ser «un Chen Tianqi con más fracasos, capaz de conservar esa mentalidad original para enfrentar lo que queremos enfrentar». Chen Tianqi estuvo de acuerdo con gusto: «Habrá más fracasos, y podré seguir conservando esa mentalidad original para enfrentarlos. Así sería lo mejor.»

\begin{importantbox}{Frases clave de la entrevista}
A continuación se presentan algunas de las frases originales más reveladoras de esta entrevista:
\begin{itemize}
    \item «Trabajar en un problema bastante mediocre y en uno importante toma casi el mismo tiempo.»
    \item «El tiempo es limitado, así que hay que hacer lo que uno realmente quiere hacer.»
    \item «Muchas veces tu valor es, en realidad, un compromiso que tu yo pasado hizo contigo mismo.»
    \item «Si eliges este camino y fracasas, ¿te arrepentirías?»
    \item «El fracaso tampoco es tan terrible. Hay que aceptar el hecho de que se puede fracasar.»
    \item «Creo que lo más difícil para mí ahora sigue siendo mantener la intención original.»
    \item «Hacer tecnología excelente siempre es acertado.»
    \item ``We need to reinvent ourselves, then keep doing the next thing.''
\end{itemize}
\end{importantbox}

\subsection{Lecturas adicionales}

\begin{itemize}
    \item Chen Tianqi, «Diez años de investigación en aprendizaje automático», Zhihu, 2019. Este artículo repasa su trayectoria de investigación de 2009 a 2019, y muchas de las historias que aparecen ahí se desarrollan con más detalle en esta entrevista. La frase del artículo «un proceso aleatorio siempre converge a un estado estable proporcional al esfuerzo invertido» se difundió ampliamente.
    \item Artículo de XGBoost: Chen \& Guestrin, ``XGBoost: A Scalable Tree Boosting System'', KDD 2016. Este artículo tiene hasta la fecha casi 70,000 citas y es uno de los artículos más citados en la historia del KDD.
    \item Artículo de MXNet: Chen et al., ``MXNet: A Flexible and Efficient Machine Learning Library for Heterogeneous Distributed Systems'', NeurIPS Workshop 2015. Los autores aparecen en orden alfabético, lo cual refleja el espíritu de «colaboración pura» entre estudiantes de doctorado.
    \item Artículo de TVM: Chen et al., ``TVM: An Automated End-to-End Optimizing Compiler for Deep Learning'', OSDI 2018. Dio origen a la línea de investigación de la compilación de aprendizaje automático.
    \item Artículo de LIME: Ribeiro, Singh, Guestrin, ``Why Should I Trust You? Explaining the Predictions of Any Classifier'', KDD 2016. Otro trabajo pionero del laboratorio de Carlos del mismo año, que abrió la línea de la interpretabilidad en el aprendizaje automático.
    \item Proyecto MLC LLM: \url{https://github.com/mlc-ai/mlc-llm}. Compila y despliega modelos de lenguaje grandes en distintos dispositivos de borde.
    \item Proyecto Apache TVM: \url{https://tvm.apache.org/}. Infraestructura de código abierto para la compilación de aprendizaje automático.
    \item Curso MLC (el curso de compilación de aprendizaje automático que Chen Tianqi imparte en CMU): \url{https://mlc.ai/}. Explica de manera sistemática los principios y la práctica de la compilación de aprendizaje automático.
    \item XGrammar: \url{https://github.com/mlc-ai/xgrammar}. Biblioteca de código abierto para la generación de datos estructurados, que se ha convertido en uno de los estándares de facto para la salida estructurada de agentes de LLM.
    \item Segundo episodio del podcast WhynotTV (entrevista a Hu Yuanming (胡渊鸣)): también aborda temas como la comercialización de proyectos de código abierto y el emprendimiento académico, y ofrece un contraste interesante con este episodio.
\end{itemize}

\end{document}
